\section*{Sets and Indices}

No nontrivial index sets are used in the implemented model. The problem has two product-specific continuous production decisions:
\[
\text{liquid fertilizer}, \qquad \text{solid fertilizer}.
\]

\section*{Parameters}

All time coefficients are measured in minutes per lot, and machine availability is converted from hours to minutes in the code.

\[
a_{1L} = 50
\]
Machine 1 time required per lot of liquid fertilizer (code: \texttt{m1\_liquid}).

\[
a_{2L} = 30
\]
Machine 2 time required per lot of liquid fertilizer (code: \texttt{m2\_liquid}).

\[
a_{1S} = 24
\]
Machine 1 time required per lot of solid fertilizer (code: \texttt{m1\_solid}).

\[
a_{2S} = 33
\]
Machine 2 time required per lot of solid fertilizer (code: \texttt{m2\_solid}).

\[
I_L^0 = 30, \qquad I_S^0 = 90
\]
Initial inventories of liquid and solid fertilizer, respectively (code: \texttt{init\_liquid}, \texttt{init\_solid}).

\[
T_1 = 40 \times 60 = 2400, \qquad T_2 = 35 \times 60 = 2100
\]
Available processing time on Machines 1 and 2 during the week, respectively, in minutes (code: \texttt{m1\_avail}, \texttt{m2\_avail}).

\[
D_L = 75, \qquad D_S = 95
\]
Forecast demand for liquid and solid fertilizer, respectively (code: \texttt{demand\_liquid}, \texttt{demand\_solid}).

\section*{Decision Variables}

\[
x_L \ge 0
\]
Lots of liquid fertilizer produced during the week (code: \texttt{x\_l}).

\[
x_S \ge 0
\]
Lots of solid fertilizer produced during the week (code: \texttt{x\_s}).

For convenience, the code also forms ending inventories as affine expressions:
\[
E_L = I_L^0 + x_L - D_L
\]
Ending inventory of liquid fertilizer (code expression: \texttt{end\_liquid}).

\[
E_S = I_S^0 + x_S - D_S
\]
Ending inventory of solid fertilizer (code expression: \texttt{end\_solid}).

\section*{Objective}

Maximize total ending inventory:
\[
\max \; E_L + E_S
\]
equivalently,
\[
\max \; \left(I_L^0 + x_L - D_L\right) + \left(I_S^0 + x_S - D_S\right).
\]

With the given data, this is
\[
\max \; (30 + x_L - 75) + (90 + x_S - 95).
\]

\section*{Constraints}

Machine 1 weekly capacity:
\[
a_{1L}x_L + a_{1S}x_S \le T_1,
\]
i.e.,
\[
50x_L + 24x_S \le 2400.
\]

Machine 2 weekly capacity:
\[
a_{2L}x_L + a_{2S}x_S \le T_2,
\]
i.e.,
\[
30x_L + 33x_S \le 2100.
\]

Liquid demand satisfaction, implemented as nonnegative ending inventory:
\[
E_L \ge 0,
\]
equivalently,
\[
I_L^0 + x_L - D_L \ge 0,
\]
or
\[
x_L \ge D_L - I_L^0 = 45.
\]

Solid demand satisfaction, implemented as nonnegative ending inventory:
\[
E_S \ge 0,
\]
equivalently,
\[
I_S^0 + x_S - D_S \ge 0,
\]
or
\[
x_S \ge D_S - I_S^0 = 5.
\]

Nonnegativity:
\[
x_L \ge 0,\qquad x_S \ge 0.
\]

\section*{Notes on Implementation}

The implemented model is a continuous linear program with two decision variables. Fractional lots are allowed because both variables are created as continuous nonnegative variables with no integrality restrictions.

The objective in the code is written using ending inventory expressions directly:
\[
\max \; (I_L^0 + x_L - D_L) + (I_S^0 + x_S - D_S).
\]
Since \(I_L^0\), \(I_S^0\), \(D_L\), and \(D_S\) are constants, this is equivalent to maximizing \(x_L + x_S\) plus a constant; however, the formulation above matches the code exactly.

The model is solved using PuLP with the Gurobi command-line solver interface (\texttt{GUROBI\_CMD}). No exact enumeration or dynamic programming is used.
